\chapter{Algorithms and Data Structures}\label{ch:iv:algorithms-and-data-structures}

``Given a stream of trades, report after each one the median size of the last thousand.'' The candidate sorts the window
every time. It works, and the interviewer asks what happens at a million trades a second. Two heaps, an order-statistic
tree, and a counting array for bounded sizes are the three answers that score, each for a different reason, and the
candidate who names the families before choosing has already done half the work. Coding interviews at trading firms ask
the standard families of algorithm questions, often dressed in market data: order books, tapes, windows of prices, halts
and intervals. This chapter lists the families, the method, and the trading-flavoured variants.

\section{The families and how to recognise them}

\begin{center}
\footnotesize
\begin{tabular}{@{}lll@{}}
\toprule
Family & The clue in the question & Typical cost \\
\midrule
Hashing & ``have we seen'', pairs summing to a target, counts & $O(n)$ \\
Two pointers & two sorted sequences, or a condition on pairs & $O(n)$ \\
Sliding window & ``the last $k$'', a window moving over a stream & $O(n)$ or $O(n \log k)$ \\
Heaps & top $k$, a running median, the next event & $O(n \log k)$ \\
\omterm{def:iv:algorithms-and-data-structures:families}{Sweep line} & intervals, ``at the same time'', overlaps & $O(n \log n)$ \\
Binary search & a sorted array, a monotone predicate & $O(\log n)$ per query \\
Graphs & conversions, dependencies, shortest paths, cycles & $O(V E)$ or $O(E \log V)$ \\
Dynamic programming & optimal choices with overlapping subproblems & states $\times$ transitions \\
\bottomrule
\end{tabular}
\end{center}

\begin{definition}[Two-pointer technique, monotonic deque, sweep line]\label{def:iv:algorithms-and-data-structures:families}
The \emph{two-pointer technique}\index{two-pointer technique} walks two indices through one or two sorted sequences, advancing
whichever the condition dictates, so that each element is visited once. A \emph{monotonic deque}\index{monotonic deque} is a
double-ended queue of indices whose values are kept in decreasing (or increasing) order, so that the maximum (or minimum) of a
sliding window is at its front at every step. A \emph{sweep line}\index{sweep line} solves interval problems by sorting the
intervals' endpoints as events and processing them in order while maintaining a running state.
\end{definition}

\omcode{code/interviews/21-algorithms-and-data-structures/python/iv_algos.py}{55}{66}{The sliding-window maximum with a
monotonic deque: every index enters and leaves the deque once, so the whole pass is linear.}\label{lst:iv:algorithms-and-data-structures:slidingmax}

\section{Trading-flavoured variants}

Market data turns the families into familiar questions. A tape is a sorted sequence (two pointers, binary search by time);
a book side is a set of price levels with add, cancel and best-price queries (a balanced tree or a heap with lazy deletion);
open orders, halts and sessions are intervals (a \omterm{def:iv:algorithms-and-data-structures:families}{sweep line}); rolling statistics are windows (a deque, two heaps, or
Welford-style updates, One Quant Book~4, chapter~25); currencies and conversion rates are a graph whose cycles with a product
above one are arbitrages (Bellman--Ford on $-\log$ of the rates).

\omcode{code/interviews/21-algorithms-and-data-structures/python/iv_algos.py}{126}{133}{A sweep line: the maximum number of
orders open at once. Sorting puts an end before a start at the same time, so an order ending when another begins does not
overlap it.}\label{lst:iv:algorithms-and-data-structures:sweep}

The book side is the one data structure every trading interview asks about, because every trading system has one (One Quant
Book~10, chapter~1, and Book~13, chapter~19). In an interview the price-level structure below is enough; the design questions
of \cref{ch:iv:system-design} ask what a production book adds.

\omcode{code/interviews/21-algorithms-and-data-structures/cpp/iv_algos.hpp}{65}{83}{One side of a book aggregated by price in
C++20: an ordered map gives logarithmic add, cancel and best price.}\label{lst:iv:algorithms-and-data-structures:levels}

\section{The method: clarify, brute force, improve, test}

\begin{method}[Solving a coding question]\label{met:iv:algorithms-and-data-structures:method}
\begin{enumerate}
\item \emph{Clarify}: input sizes and ranges, duplicates, ties, empty input, what to return when there is no answer.
\item \emph{Brute force}: say the obvious algorithm and its cost; it is the baseline and, later, the oracle.
\item \emph{Improve}: name the family, derive the better complexity, and say the lower bound (every element must be read).
\item \emph{Code}: clean names, no premature cleverness, edge cases handled.
\item \emph{Test}: run a small case by hand, then the edge cases; say how you would test at scale (next section).
\end{enumerate}
\end{method}

\section{Tutorial: testing an answer against a brute-force oracle}\label{tut:iv:algorithms-and-data-structures:oracle}

\begin{tutorial}
\textbf{Goal.} Test the sliding-window maximum of \cref{lst:iv:algorithms-and-data-structures:slidingmax} against an answer
that is obviously right, on hundreds of random inputs. \textbf{End state:} a green test that would have caught every bug
tried while writing this chapter.
\begin{tutsteps}
\item \textbf{The oracle.} Write the slowest correct answer: the maximum of every window by slicing.
\omcode{code/interviews/21-algorithms-and-data-structures/python/iv_oracles.py}{14}{15}{The oracle: quadratic and obviously
correct.}\label{lst:iv:algorithms-and-data-structures:oracle}
\item \textbf{Random small inputs.} Draw lengths from 1 to 30 and values from a small range, so that ties and repeats are
common, with a fixed seed per case.
\omcode{code/interviews/21-algorithms-and-data-structures/tests/test_solutions.py}{67}{73}{Three hundred seeded cases,
compared element by element.}\label{lst:iv:algorithms-and-data-structures:property}
\item \textbf{When it fails,} shrink the failing case by hand: remove elements while it still fails; the smallest failing
input usually names the bug.
\end{tutsteps}
The same pattern tests every answer in the chapter's code, in Python and in the C++20 twins of five of them.
\end{tutorial}

\section{Worked answers}

\begin{example}[A window with at most two venues]\label{ex:iv:algorithms-and-data-structures:venues}
\emph{``A parent order's fills are listed in time order with the venue of each. Find the longest run of consecutive fills
that touches at most two venues.''} Clarify: runs are contiguous; venues are labels. Brute force checks every window, $O(n^2)$
with a set per start. Improve with the \omterm{def:iv:algorithms-and-data-structures:families}{two-pointer technique} (\cref{def:iv:algorithms-and-data-structures:families}): extend
the right end one fill at a time, keep a count per venue in the window, and while more than two venues have a positive count,
advance the left end and decrement. Each fill enters and leaves the window once: $O(n)$ time, $O(1)$ extra space for a fixed
bound on venues. Trace on the fills A, B, A, C, C, C, B, A: after three fills the window is A, B, A; the C makes three venues, so the
left end drops the first A (still three) and the B, leaving A, C; two more C's give A, C, C, C, length 4; the B drops the
A, leaving C, C, C, B, length 4; the last A drops the three C's, leaving B, A. The answer is 4, which the brute force
confirms here and the chapter's test confirms against it on random inputs.
\end{example}

\section{Question bank}

\begin{interviewq}[$\star$ \iqroles{developer} \iqfirm{market maker}]\label{iq:iv:algorithms-and-data-structures:1}
Given a list of order sizes and a target, return two orders whose sizes sum to the target, in linear time.
\end{interviewq}

\begin{interviewq}[$\star$ \iqroles{developer} \iqfirm{proprietary firm}]\label{iq:iv:algorithms-and-data-structures:2}
Return the first client order identifier that appears twice in a stream. What if the stream does not fit in memory?
\end{interviewq}

\begin{interviewq}[$\star$ \iqroles{developer, researcher} \iqfirm{any}]\label{iq:iv:algorithms-and-data-structures:3}
Merge two trade tapes, each sorted by timestamp, into one sorted tape. What do you do with equal timestamps?
\end{interviewq}

\begin{interviewq}[$\star$ \iqroles{developer, trader} \iqfirm{market maker}]\label{iq:iv:algorithms-and-data-structures:4}
Given a day's prices, find the largest profit from one purchase followed by one sale.
\end{interviewq}

\begin{interviewq}[$\star$ \iqroles{developer, mle} \iqfirm{systematic fund}]\label{iq:iv:algorithms-and-data-structures:5}
In a sorted array of trade timestamps, find the first trade at or after a given time. What does the standard library offer?
\end{interviewq}

\begin{interviewq}[$\star\star$ \iqroles{developer} \iqfirm{proprietary firm}]\label{iq:iv:algorithms-and-data-structures:6}
Report the maximum price of every window of the last $k$ trades in $O(n)$ time overall.
\end{interviewq}

\begin{interviewq}[$\star\star$ \iqroles{developer, researcher} \iqfirm{market maker}]\label{iq:iv:algorithms-and-data-structures:7}
Report, after each trade, the median size of the last $k$ trades, faster than sorting the window each time. What if sizes are
integers below 1\,000?
\end{interviewq}

\begin{interviewq}[$\star\star$ \iqroles{developer, mle} \iqfirm{crypto firm}]\label{iq:iv:algorithms-and-data-structures:8}
From a stream of (symbol, volume) trades, return the $k$ symbols with the largest total volume. What changes if the stream is
unbounded and memory is not?
\end{interviewq}

\begin{interviewq}[$\star\star$ \iqroles{developer} \iqfirm{market maker}]\label{iq:iv:algorithms-and-data-structures:9}
Given each order's start and end time, find the largest number of orders open at the same moment.
\end{interviewq}

\begin{interviewq}[$\star\star$ \iqroles{developer, researcher} \iqfirm{any}]\label{iq:iv:algorithms-and-data-structures:10}
Merge a list of trading halts, given as time intervals, into disjoint intervals.
\end{interviewq}

\begin{interviewq}[$\star\star\star$ \iqroles{developer, trader} \iqfirm{crypto firm}]\label{iq:iv:algorithms-and-data-structures:11}
Given a matrix of conversion rates between currencies, decide whether some cycle of conversions ends with more money than it
started with.
\end{interviewq}

\begin{interviewq}[$\star\star\star$ \iqroles{developer, researcher} \iqfirm{proprietary firm}]\label{iq:iv:algorithms-and-data-structures:12}
Given a day's prices, find the largest profit with at most $k$ non-overlapping round trips, each paying a fixed fee.
\end{interviewq}

\begin{interviewq}[$\star\star\star$ \iqroles{developer} \iqfirm{market maker}]\label{iq:iv:algorithms-and-data-structures:13}
Design the data structure for one side of an order book aggregated by price, with add, cancel and best-price queries. Give the
complexities, and say what changes if prices are integers in a known narrow range.
\end{interviewq}

\begin{interviewq}[$\star\star\star$ \iqroles{developer, mle} \iqfirm{systematic fund}]\label{iq:iv:algorithms-and-data-structures:14}
Maintain the mean and variance of the last $k$ values of a stream in $O(1)$ per update, numerically stably.
\end{interviewq}

\begin{omsources}
\item T.~H.~Cormen, C.~E.~Leiserson, R.~L.~Rivest and C.~Stein, \emph{Introduction to Algorithms}, 4th edition, MIT Press,
2022.
\item B.~P.~Welford, ``Note on a method for calculating corrected sums of squares and products'', \emph{Technometrics} 4(3),
1962, 419--420.
\item One Quant Book~4, chapter~25; One Quant Book~10, chapter~1; One Quant Book~13, chapter~19.
\end{omsources}
